% --- Archon physics formalization source begin ---
% archon:physics
% archon:covers PhyXMiniProblems/problem_phyx_mini_0871.lean
% archon:source-report reports/phyx_mini/problem_phyx_mini_0871.source.json

\chapter{Physics problem phyx\_mini\_0871}
\label{ch:PhyXMiniProblems_problem_phyx_mini_0871}

\paragraph{Problem source.}
\#\# Question

What is the magnitude of the electric field at the dot?

\#\# Answer choices

- A: A: 10kV/m

- B: B: 15kV/m

- C: C: 25kV/m

- D: D: 20kV/m

\#\# Figure caption (auxiliary; use the image as primary evidence)

The image is a 2D coordinate plane with axes labeled \textbackslash{}(x\textbackslash{}) and \textbackslash{}(y\textbackslash{}). There is a right triangle formed along these axes with each leg measuring 1 cm, and an angle of 45° at the intersection. Dashed green lines run diagonally, parallel to the hypotenuse of the triangle.

The diagonal lines represent equipotential lines, labeled with electric potentials:

- One line is labeled \textbackslash{}(200 \textbackslash{}, \textbackslash{}text\{V\}\textbackslash{}).
- One line is labeled \textbackslash{}(0 \textbackslash{}, \textbackslash{}text\{V\}\textbackslash{}).
- One line is labeled \textbackslash{}(-200 \textbackslash{}, \textbackslash{}text\{V\}\textbackslash{}).

There is a black dot located centrally among these lines, indicating a specific point in the coordinate system.

\#\# Dataset metadata

Electrostatics; Physical Model Grounding Reasoning, Spatial Relation Reasoning

\paragraph{Recorded answer/context.}
D: D: 20kV/m

\paragraph{Figure/image path.}
/root/proposal\_for\_physic/science-mango/phyx\_mini\_run/phyx\_data/test\_image/871.png

\paragraph{Formalization target.}
create a compiling Lean file with sorry bodies at `PhyXMiniProblems/problem\_phyx\_mini\_0871.lean`. The Lean declarations must preserve the physical quantities, dimensions or dimensional roles, figure labels, governing-law hypotheses, and final relation expressed by this problem.
Use Mathlib/Physlib names found through LeanExplore where available. If a physics API is missing, introduce faithful local abstractions rather than scalar placeholder aliases.

\begin{theorem}[Physics formalization target]
\label{thm:physics:phyx_mini_0871:target}
\lean{PhyXMiniProblems.ProblemPhyXMini0871.electricFieldMagnitudeAtDot_matches_answer_D}
\uses{def:physics:phyx-mini-0871:phyxminiproblems-problemphyxmini0871-equipotentialfieldsetup, def:physics:phyx-mini-0871:phyxminiproblems-problemphyxmini0871-matchesprimaryequipotentialfigure, def:physics:phyx-mini-0871:phyxminiproblems-problemphyxmini0871-hasphysicalequipotentialgeometry, def:physics:phyx-mini-0871:phyxminiproblems-problemphyxmini0871-satisfiesuniformequipotentialelectrostatics, def:physics:phyx-mini-0871:phyxminiproblems-problemphyxmini0871-electricfieldmagnitudeatdotinkilovoltspermeter, def:physics:phyx-mini-0871:phyxminiproblems-problemphyxmini0871-displayedfieldmagnitudeinkilovoltspermeter, def:physics:phyx-mini-0871:phyxminiproblems-problemphyxmini0871-isuniqueexactdisplayedmatch}
The assigned autoformalize agent should translate this physics problem into Lean declarations in the covered file, with theorem and lemma proof bodies written as `by sorry`.
\end{theorem}
\begin{proof}
This is an autoformalization task, not a proof task. Produce statements that can later be proved by the physics prover without weakening the physical model.
\end{proof}

% --- Archon declaration topology begin ---
\section{Declaration topology}
\begin{definition}[electric Potential Dimension]
  \label{def:physics:phyx-mini-0871:phyxminiproblems-problemphyxmini0871-electricpotentialdimension}
  \lean{PhyXMiniProblems.ProblemPhyXMini0871.electricPotentialDimension}
  Electric potential has dimension energy per charge
\end{definition}
\begin{proof}
  This is the declared model definition of electric Potential Dimension.
\end{proof}
\begin{definition}[electric Field Dimension]
  \label{def:physics:phyx-mini-0871:phyxminiproblems-problemphyxmini0871-electricfielddimension}
  \lean{PhyXMiniProblems.ProblemPhyXMini0871.electricFieldDimension}
  Electric field has dimension potential per length
\end{definition}
\begin{proof}
  This is the declared model definition of electric Field Dimension.
\end{proof}
\begin{definition}[Length Quantity]
  \label{def:physics:phyx-mini-0871:phyxminiproblems-problemphyxmini0871-lengthquantity}
  \lean{PhyXMiniProblems.ProblemPhyXMini0871.LengthQuantity}
  A nonnegative unit-independent physical length
\end{definition}
\begin{proof}
  This is the declared model definition of Length Quantity.
\end{proof}
\begin{definition}[Electric Potential Quantity]
  \label{def:physics:phyx-mini-0871:phyxminiproblems-problemphyxmini0871-electricpotentialquantity}
  \lean{PhyXMiniProblems.ProblemPhyXMini0871.ElectricPotentialQuantity}
  \uses{def:physics:phyx-mini-0871:phyxminiproblems-problemphyxmini0871-electricpotentialdimension}
  A signed unit-independent electrostatic potential
\end{definition}
\begin{proof}
  This is the declared model definition of Electric Potential Quantity.
\end{proof}
\begin{definition}[Planar Position Quantity]
  \label{def:physics:phyx-mini-0871:phyxminiproblems-problemphyxmini0871-planarpositionquantity}
  \lean{PhyXMiniProblems.ProblemPhyXMini0871.PlanarPositionQuantity}
  A unit-independent position in the pictured two-dimensional plane
\end{definition}
\begin{proof}
  This is the declared model definition of Planar Position Quantity.
\end{proof}
\begin{definition}[Electric Field Vector Quantity]
  \label{def:physics:phyx-mini-0871:phyxminiproblems-problemphyxmini0871-electricfieldvectorquantity}
  \lean{PhyXMiniProblems.ProblemPhyXMini0871.ElectricFieldVectorQuantity}
  \uses{def:physics:phyx-mini-0871:phyxminiproblems-problemphyxmini0871-electricfielddimension}
  A unit-independent planar electric-field vector
\end{definition}
\begin{proof}
  This is the declared model definition of Electric Field Vector Quantity.
\end{proof}
\begin{definition}[x Coordinate]
  \label{def:physics:phyx-mini-0871:phyxminiproblems-problemphyxmini0871-xcoordinate}
  \lean{PhyXMiniProblems.ProblemPhyXMini0871.xCoordinate}
  Coordinate 0, the horizontal x direction in the image
\end{definition}
\begin{proof}
  This is the declared model definition of x Coordinate.
\end{proof}
\begin{definition}[y Coordinate]
  \label{def:physics:phyx-mini-0871:phyxminiproblems-problemphyxmini0871-ycoordinate}
  \lean{PhyXMiniProblems.ProblemPhyXMini0871.yCoordinate}
  Coordinate 1, the vertical y direction in the image
\end{definition}
\begin{proof}
  This is the declared model definition of y Coordinate.
\end{proof}
\begin{definition}[length In Meters]
  \label{def:physics:phyx-mini-0871:phyxminiproblems-problemphyxmini0871-lengthinmeters}
  \lean{PhyXMiniProblems.ProblemPhyXMini0871.lengthInMeters}
  \uses{def:physics:phyx-mini-0871:phyxminiproblems-problemphyxmini0871-lengthquantity}
  Read a physical length in coherent-SI metres
\end{definition}
\begin{proof}
  This is the declared model definition of length In Meters.
\end{proof}
\begin{definition}[electric Potential In Volts]
  \label{def:physics:phyx-mini-0871:phyxminiproblems-problemphyxmini0871-electricpotentialinvolts}
  \lean{PhyXMiniProblems.ProblemPhyXMini0871.electricPotentialInVolts}
  \uses{def:physics:phyx-mini-0871:phyxminiproblems-problemphyxmini0871-electricpotentialquantity}
  Read an electrostatic potential in coherent-SI volts
\end{definition}
\begin{proof}
  This is the declared model definition of electric Potential In Volts.
\end{proof}
\begin{definition}[position Vector In Meters]
  \label{def:physics:phyx-mini-0871:phyxminiproblems-problemphyxmini0871-positionvectorinmeters}
  \lean{PhyXMiniProblems.ProblemPhyXMini0871.positionVectorInMeters}
  \uses{def:physics:phyx-mini-0871:phyxminiproblems-problemphyxmini0871-planarpositionquantity}
  Read a planar physical position as a vector of metre coordinates
\end{definition}
\begin{proof}
  This is the declared model definition of position Vector In Meters.
\end{proof}
\begin{definition}[position In Physlib Space]
  \label{def:physics:phyx-mini-0871:phyxminiproblems-problemphyxmini0871-positioninphyslibspace}
  \lean{PhyXMiniProblems.ProblemPhyXMini0871.positionInPhyslibSpace}
  \uses{def:physics:phyx-mini-0871:phyxminiproblems-problemphyxmini0871-planarpositionquantity, def:physics:phyx-mini-0871:phyxminiproblems-problemphyxmini0871-positionvectorinmeters}
  Regard the coherent-SI coordinates as a point of Physlib's planar space
\end{definition}
\begin{proof}
  This is the declared model definition of position In Physlib Space.
\end{proof}
\begin{definition}[electric Field Vector In Volts Per Meter]
  \label{def:physics:phyx-mini-0871:phyxminiproblems-problemphyxmini0871-electricfieldvectorinvoltspermeter}
  \lean{PhyXMiniProblems.ProblemPhyXMini0871.electricFieldVectorInVoltsPerMeter}
  \uses{def:physics:phyx-mini-0871:phyxminiproblems-problemphyxmini0871-electricfieldvectorquantity}
  Read an electric-field vector in coherent-SI volts per metre
\end{definition}
\begin{proof}
  This is the declared model definition of electric Field Vector In Volts Per Meter.
\end{proof}
\begin{definition}[Potential Line]
  \label{def:physics:phyx-mini-0871:phyxminiproblems-problemphyxmini0871-potentialline}
  \lean{PhyXMiniProblems.ProblemPhyXMini0871.PotentialLine}
  The three dashed equipotential lines, ordered by increasing potential
\end{definition}
\begin{proof}
  This is the declared model definition of Potential Line.
\end{proof}
\begin{definition}[Adjacent Line Gap]
  \label{def:physics:phyx-mini-0871:phyxminiproblems-problemphyxmini0871-adjacentlinegap}
  \lean{PhyXMiniProblems.ProblemPhyXMini0871.AdjacentLineGap}
  The two adjacent perpendicular gaps measured in the image
\end{definition}
\begin{proof}
  This is the declared model definition of Adjacent Line Gap.
\end{proof}
\begin{definition}[Coordinate Axis]
  \label{def:physics:phyx-mini-0871:phyxminiproblems-problemphyxmini0871-coordinateaxis}
  \lean{PhyXMiniProblems.ProblemPhyXMini0871.CoordinateAxis}
  The two coordinate axes printed in the primary image
\end{definition}
\begin{proof}
  This is the declared model definition of Coordinate Axis.
\end{proof}
\begin{definition}[Figure Color]
  \label{def:physics:phyx-mini-0871:phyxminiproblems-problemphyxmini0871-figurecolor}
  \lean{PhyXMiniProblems.ProblemPhyXMini0871.FigureColor}
  Colors needed to transcribe the dashed-line styling
\end{definition}
\begin{proof}
  This is the declared model definition of Figure Color.
\end{proof}
\begin{definition}[Equipotential Line Figure]
  \label{def:physics:phyx-mini-0871:phyxminiproblems-problemphyxmini0871-equipotentiallinefigure}
  \lean{PhyXMiniProblems.ProblemPhyXMini0871.EquipotentialLineFigure}
  \uses{def:physics:phyx-mini-0871:phyxminiproblems-problemphyxmini0871-potentialline, def:physics:phyx-mini-0871:phyxminiproblems-problemphyxmini0871-adjacentlinegap, def:physics:phyx-mini-0871:phyxminiproblems-problemphyxmini0871-coordinateaxis, def:physics:phyx-mini-0871:phyxminiproblems-problemphyxmini0871-figurecolor}
  Literal qualitative and textual information visible in image 871.png. Numerical physical readouts are connected to the independent quantities in MatchesPrimaryEquipotentialFigure below
\end{definition}
\begin{proof}
  This is the declared model definition of Equipotential Line Figure.
\end{proof}
\begin{definition}[Equipotential Field Setup]
  \label{def:physics:phyx-mini-0871:phyxminiproblems-problemphyxmini0871-equipotentialfieldsetup}
  \lean{PhyXMiniProblems.ProblemPhyXMini0871.EquipotentialFieldSetup}
  \uses{def:physics:phyx-mini-0871:phyxminiproblems-problemphyxmini0871-lengthquantity, def:physics:phyx-mini-0871:phyxminiproblems-problemphyxmini0871-electricpotentialquantity, def:physics:phyx-mini-0871:phyxminiproblems-problemphyxmini0871-planarpositionquantity, def:physics:phyx-mini-0871:phyxminiproblems-problemphyxmini0871-electricfieldvectorquantity, def:physics:phyx-mini-0871:phyxminiproblems-problemphyxmini0871-potentialline, def:physics:phyx-mini-0871:phyxminiproblems-problemphyxmini0871-equipotentiallinefigure}
  Independent physical observables and geometry. The potential and field are not defined from the answer table; the electrostatic laws constrain them separately
\end{definition}
\begin{proof}
  This is the declared model definition of Equipotential Field Setup.
\end{proof}
\begin{definition}[Matches Primary Equipotential Figure]
  \label{def:physics:phyx-mini-0871:phyxminiproblems-problemphyxmini0871-matchesprimaryequipotentialfigure}
  \lean{PhyXMiniProblems.ProblemPhyXMini0871.MatchesPrimaryEquipotentialFigure}
  \uses{def:physics:phyx-mini-0871:phyxminiproblems-problemphyxmini0871-xcoordinate, def:physics:phyx-mini-0871:phyxminiproblems-problemphyxmini0871-ycoordinate, def:physics:phyx-mini-0871:phyxminiproblems-problemphyxmini0871-lengthinmeters, def:physics:phyx-mini-0871:phyxminiproblems-problemphyxmini0871-electricpotentialinvolts, def:physics:phyx-mini-0871:phyxminiproblems-problemphyxmini0871-positionvectorinmeters, def:physics:phyx-mini-0871:phyxminiproblems-problemphyxmini0871-equipotentialfieldsetup}
  All quantitative and qualitative evidence extracted from the primary image. The chosen representative points lie on a common normal through the marked dot; parallelism makes their adjacent separations equal to the displayed perpendicular spacing. No electric-field magnitude occurs in this premise
\end{definition}
\begin{proof}
  This is the declared model definition of Matches Primary Equipotential Figure.
\end{proof}
\begin{definition}[Has Physical Equipotential Geometry]
  \label{def:physics:phyx-mini-0871:phyxminiproblems-problemphyxmini0871-hasphysicalequipotentialgeometry}
  \lean{PhyXMiniProblems.ProblemPhyXMini0871.HasPhysicalEquipotentialGeometry}
  \uses{def:physics:phyx-mini-0871:phyxminiproblems-problemphyxmini0871-lengthinmeters, def:physics:phyx-mini-0871:phyxminiproblems-problemphyxmini0871-equipotentialfieldsetup}
  Positivity and normalization conditions for the physical geometry
\end{definition}
\begin{proof}
  This is the declared model definition of Has Physical Equipotential Geometry.
\end{proof}
\begin{definition}[Satisfies Uniform Equipotential Electrostatics]
  \label{def:physics:phyx-mini-0871:phyxminiproblems-problemphyxmini0871-satisfiesuniformequipotentialelectrostatics}
  \lean{PhyXMiniProblems.ProblemPhyXMini0871.SatisfiesUniformEquipotentialElectrostatics}
  \uses{def:physics:phyx-mini-0871:phyxminiproblems-problemphyxmini0871-electricpotentialinvolts, def:physics:phyx-mini-0871:phyxminiproblems-problemphyxmini0871-positionvectorinmeters, def:physics:phyx-mini-0871:phyxminiproblems-problemphyxmini0871-positioninphyslibspace, def:physics:phyx-mini-0871:phyxminiproblems-problemphyxmini0871-electricfieldvectorinvoltspermeter, def:physics:phyx-mini-0871:phyxminiproblems-problemphyxmini0871-equipotentialfieldsetup}
  The standard local model for equally spaced parallel equipotentials in a uniform electrostatic field: potential is constant along each shown line; the electric field is normal to those lines; potential differences are minus the field line integral, which reduces to ΔV = -E · Δr for this uniform field; and the dimensionful field readout at the dot agrees with Physlib's Electromagnetism.ElectricField at the same time and position.
\end{definition}
\begin{proof}
  This is the declared model definition of Satisfies Uniform Equipotential Electrostatics.
\end{proof}
\begin{definition}[electric Field Magnitude At Dot In Volts Per Meter]
  \label{def:physics:phyx-mini-0871:phyxminiproblems-problemphyxmini0871-electricfieldmagnitudeatdotinvoltspermeter}
  \lean{PhyXMiniProblems.ProblemPhyXMini0871.electricFieldMagnitudeAtDotInVoltsPerMeter}
  \uses{def:physics:phyx-mini-0871:phyxminiproblems-problemphyxmini0871-electricfieldvectorinvoltspermeter, def:physics:phyx-mini-0871:phyxminiproblems-problemphyxmini0871-equipotentialfieldsetup}
  Electric-field magnitude at the dot in coherent-SI volts per metre
\end{definition}
\begin{proof}
  This is the declared model definition of electric Field Magnitude At Dot In Volts Per Meter.
\end{proof}
\begin{definition}[electric Field Magnitude At Dot In Kilovolts Per Meter]
  \label{def:physics:phyx-mini-0871:phyxminiproblems-problemphyxmini0871-electricfieldmagnitudeatdotinkilovoltspermeter}
  \lean{PhyXMiniProblems.ProblemPhyXMini0871.electricFieldMagnitudeAtDotInKilovoltsPerMeter}
  \uses{def:physics:phyx-mini-0871:phyxminiproblems-problemphyxmini0871-equipotentialfieldsetup, def:physics:phyx-mini-0871:phyxminiproblems-problemphyxmini0871-electricfieldmagnitudeatdotinvoltspermeter}
  Electric-field magnitude at the dot in kilovolts per metre
\end{definition}
\begin{proof}
  This is the declared model definition of electric Field Magnitude At Dot In Kilovolts Per Meter.
\end{proof}
\begin{definition}[Answer Choice]
  \label{def:physics:phyx-mini-0871:phyxminiproblems-problemphyxmini0871-answerchoice}
  \lean{PhyXMiniProblems.ProblemPhyXMini0871.AnswerChoice}
  Labels attached to the four displayed choices
\end{definition}
\begin{proof}
  This is the declared model definition of Answer Choice.
\end{proof}
\begin{definition}[displayed Field Magnitude In Kilovolts Per Meter]
  \label{def:physics:phyx-mini-0871:phyxminiproblems-problemphyxmini0871-displayedfieldmagnitudeinkilovoltspermeter}
  \lean{PhyXMiniProblems.ProblemPhyXMini0871.displayedFieldMagnitudeInKilovoltsPerMeter}
  \uses{def:physics:phyx-mini-0871:phyxminiproblems-problemphyxmini0871-answerchoice}
  Printed field magnitude beside each answer, in kilovolts per metre
\end{definition}
\begin{proof}
  This is the declared model definition of displayed Field Magnitude In Kilovolts Per Meter.
\end{proof}
\begin{definition}[Is Unique Exact Displayed Match]
  \label{def:physics:phyx-mini-0871:phyxminiproblems-problemphyxmini0871-isuniqueexactdisplayedmatch}
  \lean{PhyXMiniProblems.ProblemPhyXMini0871.IsUniqueExactDisplayedMatch}
  \uses{def:physics:phyx-mini-0871:phyxminiproblems-problemphyxmini0871-equipotentialfieldsetup, def:physics:phyx-mini-0871:phyxminiproblems-problemphyxmini0871-electricfieldmagnitudeatdotinkilovoltspermeter, def:physics:phyx-mini-0871:phyxminiproblems-problemphyxmini0871-answerchoice, def:physics:phyx-mini-0871:phyxminiproblems-problemphyxmini0871-displayedfieldmagnitudeinkilovoltspermeter}
  A choice is the unique exact match when its displayed value equals the independently modeled field magnitude and no other displayed value does
\end{definition}
\begin{proof}
  This is the declared model definition of Is Unique Exact Displayed Match.
\end{proof}
\begin{lemma}[electric Field Magnitude At Dot from adjacent Equipotentials]
  \label{lem:physics:phyx-mini-0871:phyxminiproblems-problemphyxmini0871-electricfieldmagnitudeatdot-from-adjacentequipotentials}
  \lean{PhyXMiniProblems.ProblemPhyXMini0871.electricFieldMagnitudeAtDot_from_adjacentEquipotentials}
  \uses{def:physics:phyx-mini-0871:phyxminiproblems-problemphyxmini0871-lengthinmeters, def:physics:phyx-mini-0871:phyxminiproblems-problemphyxmini0871-electricpotentialinvolts, def:physics:phyx-mini-0871:phyxminiproblems-problemphyxmini0871-equipotentialfieldsetup, def:physics:phyx-mini-0871:phyxminiproblems-problemphyxmini0871-matchesprimaryequipotentialfigure, def:physics:phyx-mini-0871:phyxminiproblems-problemphyxmini0871-hasphysicalequipotentialgeometry, def:physics:phyx-mini-0871:phyxminiproblems-problemphyxmini0871-satisfiesuniformequipotentialelectrostatics, def:physics:phyx-mini-0871:phyxminiproblems-problemphyxmini0871-electricfieldmagnitudeatdotinvoltspermeter}
  The potential change between either adjacent pair, divided by its normal spacing, determines the electric-field magnitude. This is a derived relation, not a premise of the model
\end{lemma}
\begin{proof}
  The conclusion follows from the governing relations and figure readouts named in the statement of electric Field Magnitude At Dot from adjacent Equipotentials.
\end{proof}
% --- Archon declaration topology end ---
% --- Archon physics formalization source end ---
